%%%%%%%%%%%%%%%%%%%%%%%%%%%%%%%%%%%%%%%%%%%%%%%%%%%%%%%%%%%%%%%%%%%%%%
% Lớp 9 | Chương VI. Hàm số y = ax² (a ≠ 0). Phương trình bậc hai một ẩn | Giải bài toán bằng cách lập phương trình
% Dạng 9D6?21-3: Toán hình học (diện tích, kích thước)
%%%%%%%%%%%%%%%%%%%%%%%%%%%%%%%%%%%%%%%%%%%%%%%%%%%%%%%%%%%%%%%%%%%%%%


% nguồn: 2026-11-14_09_de_tuyen_sinh_10_tphcm_2025_2026 (nhập 14/11/2026)
\begin{bt}%[9D6V21-3]%[ma:d2600ea8]
% nguon: Đề thi chính thức Tuyển sinh 10 TP.HCM năm học 2025-2026 – Bài 4
Một khu vườn hình chữ nhật có chiều dài là $2x\text{ (m)}$ và chiều rộng là $x\text{ (m)}$ ($x > 4$). Bác Ba làm một lối đi xung quanh khu vườn rộng $2\text{ mét}$. Phần đất còn lại ở bên trong dùng để trồng hoa.
\begin{enumerate}[a)]
	\item Viết biểu thức theo $x$ biểu diễn diện tích phần đất dùng để trồng hoa và thu gọn biểu thức đó.
	\item Giả sử diện tích phần đất trồng hoa là $4800\text{ m}^2$. Tính chiều dài và chiều rộng của khu vườn.
\end{enumerate}
\loigiai{
\begin{enumerate}[a)]
	\item Vì lối đi quanh khu vườn rộng $2\text{ m}$ nên phần đất trồng hoa bên trong là hình chữ nhật có:
	\begin{itemize}
		\item Chiều rộng là: $x - 2\cdot 2 = x - 4\text{ (m)}$.
		\item Chiều dài là: $2x - 2\cdot 2 = 2x - 4\text{ (m)}$.
	\end{itemize}
	Biểu thức biểu diễn diện tích phần đất dùng để trồng hoa là:
	$S = (x - 4)(2x - 4) = 2x^2 - 4x - 8x + 16 = 2x^2 - 12x + 16\text{ (m}^2\text{)}$.
	\item Theo giả thiết, diện tích đất trồng hoa là $4800\text{ m}^2$ nên ta có phương trình:
	$2x^2 - 12x + 16 = 4800 \Leftrightarrow 2x^2 - 12x - 4784 = 0 \Leftrightarrow x^2 - 6x - 2392 = 0$.\\
	Ta có $\Delta' = (-3)^2 - 1\cdot(-2392) = 9 + 2392 = 2401 = 49^2 > 0$.\\
	Phương trình có hai nghiệm:
	$x_1 = 3 + 49 = 52$ (thoả mãn điều kiện $x > 4$),\\
	$x_2 = 3 - 49 = -46$ (loại vì $x > 4$).\\
	Với $x = 52$, ta có:
	\begin{itemize}
		\item Chiều rộng của khu vườn là: $x = 52\text{ (m)}$.
		\item Chiều dài của khu vườn là: $2x = 2\cdot 52 = 104\text{ (m)}$.
	\end{itemize}
	Vậy chiều dài của khu vườn là $104\text{ m}$ và chiều rộng là $52\text{ m}$.
\end{enumerate}
}
\end{bt}
